\section{Introduction}
\label{sec:intro}
Training modern deep neural networks requires huge computational power and storage due to the large amount of model weights and datasets.
Weight pruning and coreset selection \citep{hoefler2021sparsity,killamsetty2021grad} are two emerging paradigms to enhance training efficiency or model quality. Their key ideas are to develop proper rules to remove redundant model weights and to select the most informative samples during training, respectively. 
Notably, in scenarios with noisy data, coreset selection has been shown to improve model accuracy by filtering out less valuable samples.
It has been repeatedly reported in the literature \citep{hoefler2021sparsity,mirzasoleiman2020coresets,killamsetty2021grad} that the proposed algorithms under these two paradigms can effectively reduce training costs with negligible performance degradation. 

However, we notice that weight pruning and coreset selection are always investigated independently and their interaction is overlooked. In contrast, for the two corresponding techniques i.e., feature screening \footnote{Features in classical machine learning models always have one-to-one correspondence  with model weights.} and sample screening, in classical machine learning such as support vector machines \cite{shibagaki2016simultaneous,zhang2017scaling}, their interaction has been extensively studied and exploited to develop efficient training acceleration algorithms. Theoretical analysis shows that there exists a significant synergistic effect between feature and sample screening, that is, discarding the identified redundant features (resp. samples)  leads to a more accurate estimation of the primal (resp. dual) optimum, which in turn enhances the sample (resp. feature) screening rules in detecting irrelevant samples (resp. features). These results are derived mainly depending on the theoretical tools, e.g., Karush–Kuhn–Tucker (KKT) conditions \cite{kuhn2013nonlinear}, in convex optimization. Effective joint screening methods are proposed to achieve a synergistic effect, which enables model training with 10 million of features and samples possible on the resource-restricted devices such as laptop.  {\it Although the above optimal tools do not exist in deep learning due to the highly complicated nonconvex training problem, there seems to be no reason to assert that a similar mechanism of the synergy effect does not hold in weight pruning and coreset selection paradigms, which is the main motivation of this work. }

The first thing that can be expected is that the coreset selection training process appears to be inherently more susceptible to overfitting and thus would benefit more from the regularization effects of weight pruning.
To provide an in-depth and transparent analysis of the interaction between weight pruning and coreset selection in deep learning, and considering the complexity of deep neural networks, we first perform an in-depth investigation on polynomial interpolation tasks. 
We employ standard weight pruning and coreset selection techniques for neural networks in this experiment (Fig. \ref{fig:poly}) to ensure that the results reflect deep learning principles. 
We get the following two key observations. One is that an excessive number of samples to be interpolated, especially noisy ones, can cause the coefficients of the polynomial to be overtuned in order to fit all the samples. This significantly complicates the identification of irrelevant weights during pruning, as most existing pruning methods develop their importance metric explicitly or implicitly based on the weight magnitude. The other observation is that, under the standard metric in coreset selection, the difficulty of selection grows exponentially with the number of redundant model weights. Our initial findings expose a critical interaction between model weights and training samples: redundant samples complicate weight pruning, while irrelevant weights obscure coreset selection. 
This mutual interference implies the limitations of traditional approaches that treat pruning and coreset selection as independent processes.


To further investigate and harness the above interplay in deep learning, we develop a \textbf{S}imultaneous \textbf{W}eight \textbf{a}nd \textbf{S}ample \textbf{T}ailoring mechanism (SWaST) that alternately performs weight pruning and coreset selection to establish a synergistic effect in training.
Every \(\mathcal{R}\) epochs, SWaST performs coreset selection to identify the most representative samples. During subsequent \(\mathcal{R}\) epochs, it performs training and online pruning using this selected coreset.
We derive two variants based on the aggressiveness of the pruning strategy: SWaST-trim (which prunes only the fully connected layers) and SWaST-cut (which prunes the entire network). SWaST-trim utilizes a conservative pruning strategy, retaining more parameters and ensuring better training stability. In contrast, SWaST-cut removes more parameters for greater efficiency, but this can lead to training instability, referred to as “{\it critical double-loss}”. This instability occurs when pivotal weights and their supportive samples are mistakenly eliminated at the same time, causing an almost irreversible performance degradation. 
In contrast to classic machine learning models, this phenomenon can easily occur in deep learning due to the absence of theoretical guarantees in pruning and coreset selection. This underlying complexity elucidates why existing methods in pruning and coreset selection have traditionally been developed in isolation.
To ensure training stability, we incorporate a state preservation mechanism that captures the model's output following coreset selection. This mechanism helps preserve critical knowledge during training and pruning while allowing adaptation to the refined dataset and architecture. Extensive experiments reveal a strong synergy between pruning and coreset selection across varying prune rates and coreset sizes, delivering accuracy boosts of up to 17.83\% alongside 10\% to 90\% FLOPs reductions.

Our main contributions can be summarized as follows.
\begin{itemize}
\item We conduct a transparent analysis to explore and identify the fundamental interplay between weight pruning and coreset selection within deep learning.
\item We propose an alternative pruning and coreset selection method SWaST, designed to further investigate and  exploit the interplay to establish a synergistic effect.
\item We design a state preservation mechanism for SWaST that stabilizes training and maintains model performance under aggressive pruning rates and small coreset sizes.
\item Experiments verify the synergistic effect established by SWaST, with a significant accuracy gain of up to 17.83\%.
\end{itemize}